% (c) Nikita Lisitsa, lisyarus@gmail.com, 2026

\documentclass[10pt]{beamer}

\usepackage[T2A]{fontenc}
\usepackage[russian]{babel}
\usepackage{minted}

\usepackage[most]{tcolorbox}
\tcbuselibrary{minted}

\usepackage{graphicx}
\graphicspath{ {./images/} }

\usepackage{adjustbox}

\usepackage{color}
\usepackage{soul}

\usepackage{hyperref}

\usetheme{metropolis}

\definecolor{red}{rgb}{1,0,0}
\definecolor{codebg}{RGB}{29,35,49}
\setminted{fontsize=\footnotesize}

\makeatletter
\newcommand{\slideimage}[1]{
  \begin{figure}
    \begin{adjustbox}{width=\textwidth, totalheight=\textheight-2\baselineskip-2\baselineskip,keepaspectratio}
      \includegraphics{#1}
    \end{adjustbox}
  \end{figure}
}
\makeatother

\title{Язык программирования C++}
\subtitle{Лекция 10: конструктор копирования, оператор присваивания, const, перегрузка операторов}
\date{2026}
\author{Никита Лисица}

\setbeamertemplate{footline}[frame number]

\begin{document}

\frame{\titlepage}

\begin{frame}[fragile]
\frametitle{Вспоминаем прошлую лекцию}
\usemintedstyle{lightbulb}
\begin{tcblisting}{listing engine=minted, minted language=cpp, minted options={fontsize=\scriptsize}, listing only, colback=codebg, colframe=codebg, rounded corners}
class Array {
private:
  size_t size;
  int * data;

public:
  Array();
  Array(size_t size);
  ~Array();

  int get(size_t index);
  void set(size_t index, int value);
};
\end{tcblisting}
\end{frame}

\begin{frame}[fragile]
\frametitle{Конструктор копирования (copy constructor)}
\usemintedstyle{lightbulb}
\begin{itemize}
\item Часто полезно создать копию объекта
\pause
\item Можно сделать для этого отдельный метод:
\begin{tcblisting}{listing engine=minted, minted language=cpp, minted options={fontsize=\scriptsize}, listing only, colback=codebg, colframe=codebg, rounded corners}
class Array {
  ...
public:
  ...
  Array copy();
};
\end{tcblisting}
\pause
\item Но иногда хочется пользоваться объектами класса как обычными переменными, и использовать обычный синтаксис копирования:
\begin{tcblisting}{listing engine=minted, minted language=cpp, minted options={fontsize=\scriptsize}, listing only, colback=codebg, colframe=codebg, rounded corners}
Array a(10);
Array b = a;
\end{tcblisting}
\end{itemize}
\end{frame}

\begin{frame}[fragile]
\frametitle{Конструктор копирования (copy constructor)}
\usemintedstyle{lightbulb}
\begin{itemize}
\item По умолчанию объекты классов тоже умеют копироваться, как структуры, т.е. просто копируя все поля
\pause
\begin{tcblisting}{listing engine=minted, minted language=cpp, minted options={fontsize=\scriptsize}, listing only, colback=codebg, colframe=codebg, rounded corners}
Array a(10);
Array b = a;
\end{tcblisting}
\item Где проблема? \pause Оба объекта ссылаются на один и тот же массив с данными
\pause
\item \(\Longrightarrow\) Изменение данных в одном поменяет и в другом
\pause
\item \(\Longrightarrow\) Оба в деструкторе попробуют его удалить
\end{itemize}
\end{frame}

\begin{frame}[fragile]
\frametitle{Конструктор копирования (copy constructor)}
\usemintedstyle{lightbulb}
\begin{itemize}
\item Решение: конструктор копирования
\begin{tcblisting}{listing engine=minted, minted language=cpp, minted options={fontsize=\scriptsize}, listing only, colback=codebg, colframe=codebg, rounded corners}
class Array {
  ...
public:
  ...
  // Конструктор копирования
  Array(const Array& other) {
    size = other.size;
    data = new int [size];
    for (size_t i = 0; i < size; ++i)
      data[i] = other.data[i];
  }
};
\end{tcblisting}
\pause
\begin{tcblisting}{listing engine=minted, minted language=cpp, minted options={fontsize=\scriptsize}, listing only, colback=codebg, colframe=codebg, rounded corners}
// Вызывается конструктор от size_t
Array a(10);

// Вызывается конструктор копирования
Array b = a;
\end{tcblisting}
\end{itemize}
\end{frame}

\begin{frame}[fragile]
\frametitle{Конструктор копирования (copy constructor)}
\usemintedstyle{lightbulb}
\begin{itemize}
\usemintedstyle{tango}
\item \textbf{NB:} Передача в функцию по значению или \mintinline{cpp}|return| -- тоже вызов копирования:
\usemintedstyle{lightbulb}
\begin{tcblisting}{listing engine=minted, minted language=cpp, minted options={fontsize=\scriptsize}, listing only, colback=codebg, colframe=codebg, rounded corners}
Array makeArray() {
  ...
  // Вызов конструктора копирования для возвращаемого значения
  return a;
}
\end{tcblisting}
\pause
\begin{tcblisting}{listing engine=minted, minted language=cpp, minted options={fontsize=\scriptsize}, listing only, colback=codebg, colframe=codebg, rounded corners}
void print(Array a) {
  ...
}

Array a(10);

// Вызов конструктора копирования для аргумента функции
print(a);
\end{tcblisting}
\end{itemize}
\end{frame}

\begin{frame}[fragile]
\frametitle{Конструктор копирования (copy constructor)}
\usemintedstyle{lightbulb}
\begin{itemize}
\item Что не так в этом коде?
\begin{tcblisting}{listing engine=minted, minted language=cpp, minted options={fontsize=\scriptsize}, listing only, colback=codebg, colframe=codebg, rounded corners}
void print(Array a) {
  for (size_t i = 0; i < a.get_size(); i++) {
    printf("%d,", a.get(i));
  }
}
\end{tcblisting}
\pause
\item Функция использует массив только для чтения \(\Longrightarrow\) незачем делать копию (это не бесплатно), лучше передать по константной ссылке:
\begin{tcblisting}{listing engine=minted, minted language=cpp, minted options={fontsize=\scriptsize}, listing only, colback=codebg, colframe=codebg, rounded corners}
void print(const Array& a) {
  ...
}
\end{tcblisting}
\end{itemize}
\end{frame}

\begin{frame}[fragile]
\frametitle{Конструктор копирования (copy constructor)}
\usemintedstyle{lightbulb}
\begin{itemize}
\item Почему конструктор копирования не объявить вот так?
\begin{tcblisting}{listing engine=minted, minted language=cpp, minted options={fontsize=\scriptsize}, listing only, colback=codebg, colframe=codebg, rounded corners}
class Array {
  ...
public:
  Array(Array other);
}
\end{tcblisting}
\pause
\usemintedstyle{tango}
\item Массив \mintinline{cpp}|other| тоже надо как-то передать внутрь конструктора копирования
\pause
\item \(\Longrightarrow\) Бесконечная рекурсия :)
\end{itemize}
\end{frame}

\begin{frame}[fragile]
\frametitle{Оператор присваивания (copy assignment)}
\usemintedstyle{lightbulb}
\begin{itemize}
\item В чём разница между этими двумя строками?
\begin{tcblisting}{listing engine=minted, minted language=cpp, minted options={fontsize=\scriptsize}, listing only, colback=codebg, colframe=codebg, rounded corners}
/* 1 */ int b = a;
/* 2 */ b = a;
\end{tcblisting}
\pause
\usemintedstyle{tango}
\item В первом случае объект \textit{создаётся}, и \textit{инициализируется} копией \mintinline{cpp}|b|
\pause
\item Во втором случае объект \textit{уже создан}, и \textit{перезаписывается} копией \mintinline{cpp}|b|
\pause
\item Не очень важно для примитивных типов, но важно для классов!
\end{itemize}
\end{frame}

\begin{frame}[fragile]
\frametitle{Оператор присваивания (copy assignment)}
\usemintedstyle{lightbulb}
\begin{itemize}
\begin{tcblisting}{listing engine=minted, minted language=cpp, minted options={fontsize=\scriptsize}, listing only, colback=codebg, colframe=codebg, rounded corners}
Array a(10), b(20);

// Перезапись копированием
b = a;
\end{tcblisting}
\pause
\usemintedstyle{tango}
\item Так же как с копированием, по умолчанию просто скопируются все поля \(\Longrightarrow\) все те же проблемы, + утечка памяти массива \mintinline{cpp}|b|
\end{itemize}
\end{frame}

\begin{frame}[fragile]
\frametitle{Оператор присваивания (copy assignment)}
\usemintedstyle{lightbulb}
\begin{tcblisting}{listing engine=minted, minted language=cpp, minted options={fontsize=\scriptsize}, listing only, colback=codebg, colframe=codebg, rounded corners}
// array.h
class Array {
  ...
public:
  // Оператор присваивания
  void operator = (const Array& other);
};
\end{tcblisting}
\begin{tcblisting}{listing engine=minted, minted language=cpp, minted options={fontsize=\scriptsize}, listing only, colback=codebg, colframe=codebg, rounded corners}
// array.cpp

// Для разнообразия, реализовали в .cpp
void Array::operator = (const Array& other) {
  delete [] data;
  size = other.size;
  data = new int [size];
  for (size_t i = 0; i < size; ++i)
    data[i] = other.data[i];
}
\end{tcblisting}
\end{frame}

\begin{frame}[fragile]
\frametitle{Оператор присваивания (copy assignment)}
\usemintedstyle{lightbulb}
\begin{itemize}
\item Более идиоматичный вариант:
\begin{tcblisting}{listing engine=minted, minted language=cpp, minted options={fontsize=\scriptsize}, listing only, colback=codebg, colframe=codebg, rounded corners}
Array& Array::operator = (const Array& other) {
  if (&other == this) {
    return *this;
  }
  delete [] data;
  size = other.size;
  data = new int [size];
  for (size_t i = 0; i < size; ++i)
    data[i] = other.data[i];
  return *this;
}
\end{tcblisting}
\pause
\usemintedstyle{tango}
\item Зачем проверка \mintinline{cpp}|&other == this|? \pause Для самоприсваивания: \mintinline{cpp}|a = a;|
\pause
\item Зачем возвращать \mintinline{cpp}|*this|? \pause Например, для \textit{chained assignment}: \mintinline{cpp}|a = b = c;|
\end{itemize}
\end{frame}

\begin{frame}[fragile]
\frametitle{Явные/неявные конструкторы}
\begin{itemize}
\usemintedstyle{tango}
\item С конструктором \mintinline{cpp}|Array(size_t)| есть проблема:
\usemintedstyle{lightbulb}
\begin{tcblisting}{listing engine=minted, minted language=cpp, minted options={fontsize=\scriptsize}, listing only, colback=codebg, colframe=codebg, rounded corners}
void print(Array array) { ... }

// Эквивалентно print(Array(10))
// Происходит неявное преобразование
// size_t -> Array
print(10);

// Эквивалентно Array a(10)
Array a = 10;

Array someFunc() {
  // Тоже неявное преобразование
  return 10;
}
\end{tcblisting}
\pause
\usemintedstyle{tango}
\item Инициализация присваиванием \mintinline{cpp}|a = 10| выглядит странно
\pause
\item Вызов \mintinline{cpp}|print(10)| может делать вообще не то, что мы задумали, особенно с учётом перегрузки функций
\pause
\item \mintinline{cpp}|return 10;| -- скорее всего, опечатка
\end{itemize}
\end{frame}

\begin{frame}[fragile]
\frametitle{Явные/неявные конструкторы}
\begin{itemize}
\item Решение: конструктор можно пометить как \textbf{явный (explicit)}:
\usemintedstyle{lightbulb}
\begin{tcblisting}{listing engine=minted, minted language=cpp, minted options={fontsize=\scriptsize}, listing only, colback=codebg, colframe=codebg, rounded corners}
class Array {
public:
  explicit Array(size_t size);
  ...
};

// OK, это явный вызов конструктора
Array a(10);

// Ошибка компиляции
Array a = 10;

// Ошибка компиляции: неявное
// преобразование запрещено
print(10);

Array someFunc() {
  // Ошибка компиляции
  return 10;
}
\end{tcblisting}
\end{itemize}
\end{frame}

\begin{frame}[fragile]
\frametitle{(N)RVO, copy elision}
\begin{itemize}
\item Во что компилируется такой код?
\usemintedstyle{lightbulb}
\begin{tcblisting}{listing engine=minted, minted language=cpp, minted options={fontsize=\scriptsize}, listing only, colback=codebg, colframe=codebg, rounded corners}
Array makeArray() {
  ...
  return Array(10);
}
\end{tcblisting}
\pause
\usemintedstyle{tango}
\item Функция, вызывающая \mintinline{cpp}|makeArray()|, заводит на стеке место под возвращаемое значение, и передаёт его адрес в \mintinline{cpp}|makeArray()|
\item \mintinline{cpp}|makeArray()| создаёт локальный объект \mintinline{cpp}|Array(10)|, затем \textit{копирует} его в возвращаемое значение
\pause
\item Временный объект \mintinline{cpp}|Array(10)| после этого сразу удалится \(\Longrightarrow\) можно его создать сразу на месте возвращаемого значения
\pause
\item Эта оптимизация называется (U)RVO -- \textit{(Unnamed) Return Value Optimization}
\end{itemize}
\end{frame}

\begin{frame}[fragile]
\frametitle{(N)RVO, copy elision}
\begin{itemize}
\item Вариант чуть сложнее:
\usemintedstyle{lightbulb}
\begin{tcblisting}{listing engine=minted, minted language=cpp, minted options={fontsize=\scriptsize}, listing only, colback=codebg, colframe=codebg, rounded corners}
Array makeArray() {
  Array array(10);
  ...
  return array;
}
\end{tcblisting}
\pause
\usemintedstyle{tango}
\item Здесь объект создаётся задолго до \mintinline{cpp}|return|, тоже копируется в возвращаемое значение и потом не нужен \(\Longrightarrow\) его тоже можно создать сразу на месте возвращаемого значения
\pause
\item Эта оптимизация называется NRVO -- \textit{Named Return Value Optimization}
\end{itemize}
\end{frame}

\begin{frame}[fragile]
\frametitle{(N)RVO, copy elision}
\begin{itemize}
\item Общее название для этих оптимизаций -- \textit{copy elision}
\pause
\item Их делают компиляторы за нас :)
\pause
\item Начиная с C++17 URVO обязательна (т.е. всегда вставляется компилятором)
\pause
\item После C++11 большинство случаев copy elision свелись к \textit{implicit move}, об этом позже
\pause
\item Про это важно помнить, так как компилятор может убрать вызов копирования даже если в нём есть побочные эффекты!
\end{itemize}
\end{frame}

\begin{frame}[fragile]
\frametitle{Конструкторы полей}
\begin{itemize}
\item Если у класса есть поля, любой его конструктор должен проинициализировать все поля (т.е. вызвать их конструкторы)
\pause
\item По умолчанию, вызовется дефолтный конструктор (который для примитивных типов ничего не делает):
\usemintedstyle{lightbulb}
\begin{tcblisting}{listing engine=minted, minted language=cpp, minted options={fontsize=\scriptsize}, listing only, colback=codebg, colframe=codebg, rounded corners}
Array::Array(size_t size) {
  // Технически, к этому моменту вызвались дефолтные
  // конструкторы size и data!
  this->size = size;
  data = new int [size];
}
\end{tcblisting}
\pause
\item Что, если у поля нет дефолтного конструктора?
\end{itemize}
\end{frame}

\begin{frame}[fragile]
\frametitle{Конструкторы полей}
\begin{itemize}
\usemintedstyle{lightbulb}
\begin{tcblisting}{listing engine=minted, minted language=cpp, minted options={fontsize=\scriptsize}, listing only, colback=codebg, colframe=codebg, rounded corners}
class Vector3 {
public:
  Vector3(float x, float y, float z);
};

class RigidBody {
private:
  Vector3 position;
public:
  RigidBody() {
    // Ошибка: у поля position нет дефолтного конструктора
  }
};
\end{tcblisting}
\end{itemize}
\end{frame}

\begin{frame}[fragile]
\frametitle{Конструкторы полей}
\begin{itemize}
\item Есть специальный синтаксис для инициализации полей:
\usemintedstyle{lightbulb}
\begin{tcblisting}{listing engine=minted, minted language=cpp, minted options={fontsize=\scriptsize}, listing only, colback=codebg, colframe=codebg, rounded corners}

class RigidBody {
private:
  Vector3 position;
public:
  RigidBody() : position(0.f, 0.f, 0.f) {
    // Не ошибка: у поля position явно вызван конструктор
  }
};
\end{tcblisting}
\end{itemize}
\end{frame}

\begin{frame}[fragile]
\frametitle{Конструкторы полей}
\begin{itemize}
\item Если полей несколько, перечисляем через запятую (в порядке их объявления):
\usemintedstyle{lightbulb}
\begin{tcblisting}{listing engine=minted, minted language=cpp, minted options={fontsize=\scriptsize}, listing only, colback=codebg, colframe=codebg, rounded corners}
class RigidBody {
private:
  Vector3 position;
  Quaternion rotation;
public:
  RigidBody(float zAngle)
    : position(0.f, 0.f, 0.f)
    , rotation(cos(zAngle / 2.f), 0.f, 0.f, sin(zAngle / 2.f))
  {}
};
\end{tcblisting}
\pause
\item Если какое-то поле пропущено, у него вызывается дефолтный конструктор
\pause
\item В теле (коде) конструктора все поля класса \textbf{уже созданы и проинициализированы}!
\end{itemize}
\end{frame}

\begin{frame}[fragile]
\frametitle{Напоминание: const в языке C}
\begin{itemize}
\item Нужен, чтобы более чётко описывать интерфейсы и обезопасить себя от ошибок:
\usemintedstyle{lightbulb}
\begin{tcblisting}{listing engine=minted, minted language=cpp, minted options={fontsize=\scriptsize}, listing only, colback=codebg, colframe=codebg, rounded corners}
const int pi = 3;
// Ошибка компиляции
pi = 4;

int i = 10;
const int * ptr = &i;
// Ошибка компиляции
*ptr = 42;

// Ошибка компиляции
int * ptr2 = &pi;

// Явно говорит, что функция не меняет
// данные в переданном массиве
size_t strlen(const char* str);
\end{tcblisting}
\end{itemize}
\end{frame}

\begin{frame}[fragile]
\frametitle{Напоминание: const в языке C}
\begin{itemize}
\usemintedstyle{lightbulb}
\begin{tcblisting}{listing engine=minted, minted language=cpp, minted options={fontsize=\scriptsize}, listing only, colback=codebg, colframe=codebg, rounded corners}
// Это инициализация массива набором символов
char str1[] = "Hello, world!";

// Это инициализация указателя адресом в статической памяти
const char* str2 = "Hello, world!";

// Ошибка компиляции
char* str3 = "Hello, world!";

// Вероятнее всего, runtime креш
((char*)str2)[0] = "X";
\end{tcblisting}
\end{itemize}
\end{frame}

\begin{frame}[fragile]
\frametitle{Классы и const: константное поле}
\begin{itemize}
\item Поля класса могут быть константными, их нельзя менять:
\usemintedstyle{lightbulb}
\begin{tcblisting}{listing engine=minted, minted language=cpp, minted options={fontsize=\scriptsize}, listing only, colback=codebg, colframe=codebg, rounded corners}
class Array {
private:
  const size_t size;
  ...
public:
  void resize(size_t new_size) {
    // Ошибка компиляции: size - константное поле
    size = new_size;
  }
};
\end{tcblisting}
\end{itemize}
\end{frame}

\begin{frame}[fragile]
\frametitle{Классы и const: константное поле}
\begin{itemize}
\item Константное поле обязательно нужно явно проинициализировать в конструкторе:
\usemintedstyle{lightbulb}
\begin{tcblisting}{listing engine=minted, minted language=cpp, minted options={fontsize=\scriptsize}, listing only, colback=codebg, colframe=codebg, rounded corners}
class Array {
private:
  const size_t size;
  ...
public:
  Array() {
    // Ошибка компиляции: size не проинициализирован
  }

  Array(size_t size)
    // OK: первый size это поле, второй - аргумент конструктора
    : size(size)
  {}
};
\end{tcblisting}
\end{itemize}
\end{frame}

\begin{frame}[fragile]
\frametitle{Классы и const: константный метод}
\begin{itemize}
\item Методы тоже могут быть константными -- их можно вызывать у константных объектов класса
\pause
\item По смыслу это методы, не меняющие состояние класса
\pause
\usemintedstyle{lightbulb}
\begin{tcblisting}{listing engine=minted, minted language=cpp, minted options={fontsize=\scriptsize}, listing only, colback=codebg, colframe=codebg, rounded corners}
class Array {
private:
  size_t size;
  ...
public:
  size_t get_size() const {
    return size;
  }

  int get(size_t index) const {
    return data[index];
  }

  // Не const
  void set(size_t index, int value) {
    data[index] = value;
  }
};
\end{tcblisting}
\end{itemize}
\end{frame}

\begin{frame}[fragile]
\frametitle{Классы и const: константный метод}
\usemintedstyle{lightbulb}
\begin{itemize}
\begin{tcblisting}{listing engine=minted, minted language=cpp, minted options={fontsize=\scriptsize}, listing only, colback=codebg, colframe=codebg, rounded corners}
const Array array(10);

// OK, константный метод
array.get_size();

// OK, константный метод
array.get(5);

// Ошибка компиляции
array.set(5, 67);
\end{tcblisting}
\pause
\item Константные методы можно вызывать у любых объектов
\pause
\item Не-константные методы можно вызывать только у не-константных объектов
\end{itemize}
\end{frame}

\begin{frame}[fragile]
\frametitle{Классы и const: mutable}
\begin{itemize}
\item Рассмотрим пример кода -- данные, которые лениво (при первом обращении) читаются с диска:
\usemintedstyle{lightbulb}
\begin{tcblisting}{listing engine=minted, minted language=cpp, minted options={fontsize=\scriptsize}, listing only, colback=codebg, colframe=codebg, rounded corners}
class InputData {
private:
  Array values;
public:
  const Array& get() const {
    if (values.empty()) {
      InputFile file(...);
      values.resize(file.readNext());
      for (size_t i = 0; i < values.size(); ++i)
        values.set(i, file.readNext());
    }

    return values;
  }
}
\end{tcblisting}
\pause
\usemintedstyle{tango}
\item Ошибка компиляции: нельзя модифицировать поле \mintinline{cpp}|values| в константном методе
\end{itemize}
\end{frame}

\begin{frame}[fragile]
\frametitle{Классы и const: mutable}
\begin{itemize}
\usemintedstyle{tango}
\item Решение №1: убрать \mintinline{cpp}|const| у метода \mintinline{cpp}|get()| \pause -- не позволит передавать \mintinline{cpp}|InputData| по константной ссылке и т.п.
\pause
\item Решение №2: ключевое слово \mintinline{cpp}|mutable|
\end{itemize}
\end{frame}

\begin{frame}[fragile]
\frametitle{Классы и const: mutable}
\begin{itemize}
\usemintedstyle{tango}
\item \mintinline{cpp}|mutable| у поля класса означает, что это поле можно менять \textit{даже из константных методов} класса
\pause
\usemintedstyle{lightbulb}
\begin{tcblisting}{listing engine=minted, minted language=cpp, minted options={fontsize=\scriptsize}, listing only, colback=codebg, colframe=codebg, rounded corners}
class InputData {
private:
  mutable Array values;
public:
  const Array& get() const {
    if (values.empty()) {
      InputFile file(...);
      // OK: mutable поле можно менять
      values.resize(file.readNext());
      ...
    }

    return values;
  }
}
\end{tcblisting}
\end{itemize}
\end{frame}

\begin{frame}[fragile]
\frametitle{Классы и const: mutable}
\begin{itemize}
\usemintedstyle{tango}
\item Почти всегда считается, что использование \mintinline{cpp}|mutable| -- плохой тон, и усложняет чтение кода
\pause
\item Два основных применения, где без \mintinline{cpp}|mutable| часто не обойтись: \textit{кеширование} и \textit{мьютексы}
\pause
\item \textbf{NB:} Ещё бывают \mintinline{cpp}|mutable| лямбды, это не плохой тон, но о них позже :)
\end{itemize}
\end{frame}

\begin{frame}[fragile]
\frametitle{Операторы}
\begin{itemize}
\usemintedstyle{tango}
\item В C++ очень много операторов (\(\geq\) 60), вот некоторые из них:
\usemintedstyle{lightbulb}
\begin{tcblisting}{listing engine=minted, minted language=cpp, minted options={fontsize=\scriptsize}, listing only, colback=codebg, colframe=codebg, rounded corners}
a++ ++a a-- --a +a -a !a ~a
a() a[] a.b a->b a.*b *a &a
a+b a-b a*b a/b a%b a>>b a<<b
a&b a|b a^b a&&b a||b
a==b a!=b a<b a<=b a>b a>=b a<=>b
a=b a+=b a-=b a*=b a/=b a%=b
a&=b a|=b a^=b a>>=b a<<=b
a?b:c a,b
\end{tcblisting}
\pause
\item Как и функции, их можно перегружать, но только если один из аргументов -- пользовательский тип (т.е. класс)
\pause
\usemintedstyle{tango}
\item Перегружать можно почти всё, исключения -- оператор доступа \mintinline{cpp}|a.b|, оператор доступа по указателю на поле \mintinline{cpp}|a.*b|, тернарный условный оператор \mintinline{cpp}|a?b:c|
\end{itemize}
\end{frame}

\begin{frame}[fragile]
\frametitle{Перегрузка операторов}
\begin{itemize}
\usemintedstyle{tango}
\item Рассмотрим класс длинной арифметики \mintinline{cpp}|BigInt|: как для него определить оператор сложения? Два варианта:
\pause
\item Как свободную функцию
\usemintedstyle{lightbulb}
\begin{tcblisting}{listing engine=minted, minted language=cpp, minted options={fontsize=\scriptsize}, listing only, colback=codebg, colframe=codebg, rounded corners}
BigInt operator+(const BigInt& a, const BigInt& b) {
  ...
}
\end{tcblisting}
\pause
\item Как особый метод класса
\usemintedstyle{lightbulb}
\begin{tcblisting}{listing engine=minted, minted language=cpp, minted options={fontsize=\scriptsize}, listing only, colback=codebg, colframe=codebg, rounded corners}
BigInt BigInt::operator+(const BigInt& b) {
  // Здесь левый аргумент это this
  ...
}
\end{tcblisting}
\end{itemize}
\end{frame}

\begin{frame}[fragile]
\frametitle{Перегрузка операторов}
\begin{itemize}
\usemintedstyle{lightbulb}
\begin{tcblisting}{listing engine=minted, minted language=cpp, minted options={fontsize=\scriptsize}, listing only, colback=codebg, colframe=codebg, rounded corners}
BigInt a, b;
...
// Вызывает operator+(a, b) либо a.operator+(b)
// В зависимости от того, кто из них определён
BitInt c = a + b;
\end{tcblisting}
\end{itemize}
\end{frame}

\begin{frame}[fragile]
\frametitle{Перегрузка операторов}
\begin{itemize}
\usemintedstyle{tango}
\item Как лучше: внутри класса \mintinline{cpp}|BigInt::operator+(BigInt)| или снаружи \mintinline{cpp}|operator+(BigInt,BigInt)|?
\pause
\item Пусть есть конструктор \mintinline{cpp}|BigInt(int)|. Рассмотрим ситуации:
\usemintedstyle{lightbulb}
\begin{tcblisting}{listing engine=minted, minted language=cpp, minted options={fontsize=\scriptsize}, listing only, colback=codebg, colframe=codebg, rounded corners}
BigInt a = 10, b = 20;

// OK в любом случае
a + b;

// OK в любом случае
a + 3;

// Не OK в случае метода: operator+ объявлен методом в классе
// BigInt и не имеет отношения к типу int
3 + a;

// Метода int::operator+(BigInt) быть не может,
// так как int - не класс :)
\end{tcblisting}
\end{itemize}
\end{frame}

\begin{frame}[fragile]
\frametitle{Метод vs функция}
\begin{itemize}
\usemintedstyle{tango}
\item В целом, часто полезно придерживаться следующего правила:
\pause
\begin{itemize}
\item Если функции нужно знание о внутренностях класса, то это должен быть метод
\pause
\item Если функции достаточно публичного интерфейса класса, то лучше её сделать свободной функций, принимающей объект класса (например, по ссылке)
\end{itemize}
\pause
\item Это облегчает чтение кода класса и проверку его инвариантов
\pause
\item \textbf{NB:} Применимо не всегда, иногда удобство использования важнее чистоты кода :)
\end{itemize}
\end{frame}

\begin{frame}[fragile]
\frametitle{Compound assignment}
\begin{itemize}
\usemintedstyle{tango}
\item Операторы вида \mintinline{cpp}|a += b| называются \textit{compound assignment}
\pause
\item \mintinline{cpp}|a += b| можно реализовать через \mintinline{cpp}|a + b|, а можно и наоборот
\pause
\item Наоборот (\mintinline{cpp}|a + b| через \mintinline{cpp}|a += b|) часто оказывается эффективнее, и тогда \mintinline{cpp}|a += b| приходится делать методом класса, так как ему нужен доступ к внутренним данным класса:
\pause
\usemintedstyle{lightbulb}
\begin{tcblisting}{listing engine=minted, minted language=cpp, minted options={fontsize=\scriptsize}, listing only, colback=codebg, colframe=codebg, rounded corners}
BigInt& BigInt::operator+=(const BigInt& other) {
  // выделяем память, прибавляем, ...
  return *this;
}

BigInt operator+(BigInt a, const BigInt& b) {
  return a += b;
}
\end{tcblisting}
\end{itemize}
\end{frame}

\begin{frame}[fragile]
\frametitle{Оператор доступа по индексу}
\begin{itemize}
\usemintedstyle{lightbulb}
\begin{tcblisting}{listing engine=minted, minted language=cpp, minted options={fontsize=\scriptsize}, listing only, colback=codebg, colframe=codebg, rounded corners}
class Array {
  ...
public:

  // Можно было вернуть const int&
  int operator[](size_t index) const {
    return data[index];
  }

  int& operator[](size_t index) {
    return data[index];
  }
}
\end{tcblisting}
\pause
\item Классический паттерн в C++: приходится реализовывать два оператора -- const и не-const
\end{itemize}
\end{frame}

\begin{frame}[fragile]
\frametitle{Оператор доступа по индексу}
\begin{itemize}
\usemintedstyle{tango}
\item Начиная с C++23 есть \mintinline{cpp}|opeator[]| с произвольным числом аргументов:
\usemintedstyle{lightbulb}
\begin{tcblisting}{listing engine=minted, minted language=cpp, minted options={fontsize=\scriptsize}, listing only, colback=codebg, colframe=codebg, rounded corners}
class Matrix {
  ...
public:
  float& operator[](size_t row, size_t column);
};

Matrix m;
m[3, 4] = 2.718f;
\end{tcblisting}
\end{itemize}
\end{frame}

\begin{frame}[fragile]
\frametitle{Префиксный/постфиксный инкремент/декремент}
\begin{itemize}
\usemintedstyle{tango}
\item Напоминание:
\usemintedstyle{lightbulb}
\begin{tcblisting}{listing engine=minted, minted language=cpp, minted options={fontsize=\scriptsize}, listing only, colback=codebg, colframe=codebg, rounded corners}
int a = 10;

// Префиксный: сначала прибавить, потом вернуть значение
int b = ++a;
// Здесь a = 11, b = 11

// Постфиксный: сначала вернуть значение, потом прибавить
int c = a++;
// Здесь a = 12, c = 11
\end{tcblisting}
\pause
\item Как их перегрузить по-отдельности в C++?
\end{itemize}
\end{frame}

\begin{frame}[fragile]
\frametitle{Префиксный/постфиксный инкремент/декремент}
\begin{itemize}
\usemintedstyle{lightbulb}
\begin{tcblisting}{listing engine=minted, minted language=cpp, minted options={fontsize=\scriptsize}, listing only, colback=codebg, colframe=codebg, rounded corners}
class BigInt {
...
public:
  // Префиксный
  BigInt& operator++();

  // Постфиксный (возвращает копию!)
  // int - фикстивный параметр, не используется
  BigInt operator++(int) {
    BigInt copy(*this);
    // Или this->operator++();
    ++(*this);
    return copy;
  }
};
\end{tcblisting}
\end{itemize}
\end{frame}

\begin{frame}[fragile]
\frametitle{Операторы сравнения}
\begin{itemize}
\item В общем случае, нужно реализовать 6 штук, 4 из которых в одну строку:
\usemintedstyle{lightbulb}
\begin{tcblisting}{listing engine=minted, minted language=cpp, minted options={fontsize=\scriptsize}, listing only, colback=codebg, colframe=codebg, rounded corners}
bool operator==(const BigInt& a, const BigInt& b) { ... }
bool operator<(const BigInt& a, const BigInt& b) { ... }

bool operator!=(const BigInt& a, const BigInt& b) {
  return !(a == b);
}

bool operator>(const BigInt& a, const BigInt& b) {
  return b < a;
}

bool operator<=(const BigInt& a, const BigInt& b) {
  return !(b < a);
}

bool operator>=(const BigInt& a, const BigInt& b) {
  return !(a < b);
}
\end{tcblisting}
\end{itemize}
\end{frame}

\begin{frame}[fragile]
\frametitle{Операторы сравнения}
\begin{itemize}
\usemintedstyle{tango}
\item Можно даже \mintinline{cpp}|a == b| выразить через \mintinline{cpp}|a < b|:
\usemintedstyle{lightbulb}
\begin{tcblisting}{listing engine=minted, minted language=cpp, minted options={fontsize=\scriptsize}, listing only, colback=codebg, colframe=codebg, rounded corners}
bool operator==(const BigInt& a, const BigInt& b) {
  return !(a < b) && !(b < a);
}
\end{tcblisting}
\pause
\item Такой семантики придерживаются некоторые стандартные алгоритмы и структуры данных
\pause
\item Часто такая реализация менее эффективна, чем сравнение напрямую
\end{itemize}
\end{frame}

\begin{frame}[fragile]
\frametitle{Default и three-way сравнения}
\begin{itemize}
\usemintedstyle{tango}
\item Начиная с C++20 есть т.н. \textit{three-way comparison} (он же -- \textit{spaceship operator}): \mintinline{cpp}|a <=> b|
\pause
\item Возвращает особый тип из трёх значений (меньше, равно, больше), например \mintinline{cpp}|std::strong_ordering|
\pause
\item Можно определить только \mintinline{cpp}|a == b| и \mintinline{cpp}|a <=> b|, остальные компилятор выведет за нас
\end{itemize}
\end{frame}

\begin{frame}[fragile]
\frametitle{Оператор приведения типа}
\begin{itemize}
\usemintedstyle{tango}
\item Что, если я хочу, чтобы мой класс можно было приводить к другому типу?
\usemintedstyle{lightbulb}
\begin{tcblisting}{listing engine=minted, minted language=cpp, minted options={fontsize=\scriptsize}, listing only, colback=codebg, colframe=codebg, rounded corners}
B b;
A a = b;
A a(b);
funcOfA(b);
\end{tcblisting}
\pause
\usemintedstyle{tango}
\item Два варианта: сделать у \mintinline{cpp}|A| конструктор от \mintinline{cpp}|B|, или сделать у \mintinline{cpp}|B| оператор приведения типа к \mintinline{cpp}|A|
\pause
\item \textbf{NB:} Если \mintinline{cpp}|A| -- примитивный тип, ему не получится добавить конструктор, и остаётся только второй вариант
\end{itemize}
\end{frame}

\begin{frame}[fragile]
\frametitle{Оператор приведения типа}
\begin{itemize}
\usemintedstyle{lightbulb}
\begin{tcblisting}{listing engine=minted, minted language=cpp, minted options={fontsize=\scriptsize}, listing only, colback=codebg, colframe=codebg, rounded corners}
class BigInt {
public:
  // Оператор приведения типа BigInt -> double
  operator double() const {
    ...
  }
};

BigInt x = 10;
x = pow(x, 10000);

double y = x;
\end{tcblisting}
\end{itemize}
\end{frame}

\begin{frame}[fragile]
\frametitle{Оператор приведения типа}
\begin{itemize}
\usemintedstyle{tango}
\item Можно пометить как \mintinline{cpp}|explicit| и разрешить только явное приведение типов:
\usemintedstyle{lightbulb}
\begin{tcblisting}{listing engine=minted, minted language=cpp, minted options={fontsize=\scriptsize}, listing only, colback=codebg, colframe=codebg, rounded corners}
class BigInt {
public:
  explicit operator double() const;
};

BigInt x = 10;

// Ошибка: неявное приведение
double y = x;

// OK: явное приведение
double z = (double)x;
\end{tcblisting}
\pause
\usemintedstyle{tango}
\item Лучше всегда помечать \mintinline{cpp}|explicit|, цепочки неявных преобразований типов -- большая головная боль
\end{itemize}
\end{frame}

\begin{frame}[fragile]
\frametitle{Operator bool}
\begin{itemize}
\usemintedstyle{tango}
\item До C++11 приходилось писать странные трюки, чтобы класс можно было использовать аналогично \mintinline{cpp}|bool| -- использовать в \mintinline{cpp}|if|, \mintinline{cpp}|while|, \mintinline{cpp}|&&| и так далее (см. \textit{safe bool idiom})
\pause
\item В C++11 появились \mintinline{cpp}|explicit| операторы приведения типов, и теперь всё это делается просто как \mintinline{cpp}|explicit operator bool()|
\end{itemize}
\end{frame}

\begin{frame}[fragile]
\frametitle{Operator bool}
\usemintedstyle{lightbulb}
\begin{tcblisting}{listing engine=minted, minted language=cpp, minted options={fontsize=\scriptsize}, listing only, colback=codebg, colframe=codebg, rounded corners}
class BigInt {
public:
  explicit operator bool() const {
    return (*this) != 0;
  }
};

BigInt x = 10;

// Ошибка: неявное приведение типов
bool b = x;

// А тут можно :)
if (x) {
  ...
}
\end{tcblisting}
\end{frame}

\end{document}